\section{Literature Review / Theoretical Background}
Image focus detection aims to distinguish sharp regions from blurred regions
within an image, which is important for tasks such as image enhancement, object
recognition, and quality assessment. Early approaches mainly relied on
handcrafted image features, such as edge strength, gradient magnitude, or local
texture, combined with classical machine learning classifiers.\cite{yi2016lbp}
These methods are efficient and interpretable, but their performance depends
heavily on the quality of feature design and can degrade under varying image
conditions.

Common traditional classifiers include\cite{murphy2022pml}:

Support Vector Machines (SVMs) separate data points by finding the optimal
hyperplane that maximizes the margin between classes. They are effective in
high-dimensional spaces and small datasets. SVMs work well when classes are
linearly separable or can be transformed to be so with kernel functions.

Decision Trees classify samples by recursively splitting the data based on
feature thresholds. Each internal node represents a feature decision, and each
leaf node represents a class label. They are intuitive and easy to interpret
but can overfit if the tree grows too deep.

Random Forests are ensembles of decision trees that combine predictions from
multiple trees. They reduce overfitting and improve generalization through
majority voting (for classification) or averaging (for regression). Random
forests are robust to noisy data and can handle high-dimensional features.

Clustering Methods group similar data points based on a similarity metric,
without using labeled data. They are often used for unsupervised segmentation
or region grouping in images. Popular algorithms include K-means and
hierarchical clustering.

Convolutional neural networks (CNNs) have recently become popular in image
analysis due to their ability to automatically learn hierarchical features from
raw data. CNN-based methods can be more robust than traditional approaches, but
they usually require a large number of labeled samples and substantial
computational resources.\cite{zeng2018localmetric} Lightweight CNN
architectures reduce the number of parameters and computational cost, but it is
not clear whether they maintain an advantage over traditional methods when
training data is limited.

This study focuses on comparing traditional machine learning methods and
lightweight CNNs for segmentation-based focus detection, particularly under
low-sample conditions. The goal is to understand the trade-offs between
robustness, accuracy, and efficiency, and to identify when each approach is
most suitable.